\chapter{Limitaciones y amenazas a la validez}
\label{chap:limitaciones}

La principal limitación es la multiplicidad: se evaluaron 66 configuraciones y el Deflated Sharpe
es 0,930, inferior al umbral 0,95. Por ello el trabajo no afirma que la rentabilidad ajustada por
riesgo sobreviva concluyentemente a todas las pruebas de selección múltiple. Esta salvedad no anula
el Rank-IC, pero delimita la afirmación económica.

El horizonte de doce meses crea fuerte solapamiento entre cohortes: 117 observaciones mensuales
equivalen aproximadamente a nueve independientes. Además, la era reservada contiene sólo seis
cohortes cerradas y una observación independiente efectiva; su Rank-IC es indeterminado aunque la
confirmación económica sea favorable. El alfa factorial de la ventana de selección tampoco es
significativo por sí solo ($t=0,82$).

Existen límites de datos: universo restringido a Estados Unidos y al S\&P 500, cobertura desigual
de fundamentales, posibles restatements no observables en años antiguos y costes constantes que no
modelan impacto de mercado. Finalmente, el predominio de \textit{risk} obliga a interpretar con
prudencia la novedad de la señal; la neutralización conserva 84,35\%, lo que es evidencia parcial
contra una explicación puramente factorial, no una prueba definitiva.

Estas limitaciones conducen a una lectura disciplinada: la investigación demuestra una ordenación
predictiva robusta bajo un protocolo concreto y sugiere utilidad económica, pero requiere más
cohortes cerradas, universos adicionales y un catálogo más estrecho antes de extrapolarla.

\section{Amenazas a la validez interna}

El diseño point-in-time reduce de forma directa el \textit{lookahead bias}, pero depende de la
calidad de las fechas de filing y de la correspondencia entre métricas de Finnhub y periodos
contables. Los restatements históricos pueden ser invisibles en años tempranos: que una cifra esté
hoy disponible con una fecha de filing antigua no garantiza que sea idéntica a la cifra que conoció
el mercado ese día. La mitigación consiste en conservar periodo, filing y lag de ejecución, y en no
presentar la fuente como una cinta histórica perfecta.

La pertenencia histórica al índice limita el sesgo de supervivencia de composición, pero no elimina
la cobertura desigual de empresas retiradas ni la posible pérdida de fundamentales antiguos. El
estudio mide la fracción de filas utilizables y protege contra tickers reciclados; una réplica con
una fuente comercial point-in-time sería una prueba adicional de robustez. Esta limitación afecta
principalmente a la generalización, no a la trazabilidad del panel que sí se evaluó.

\section{Amenazas estadísticas y económicas}

Las 117 cohortes mensuales no son 117 ensayos independientes: el horizonte de doce meses hace que
dos cohortes contiguas compartan casi toda su ventana de retorno. Newey--West y bootstrap por bloques
reducen la sobreconfianza, pero no crean información nueva; por ello se informa también de unas
nueve observaciones efectivas. La multiplicidad añade una segunda cautela: las 66 configuraciones y
un Deflated Sharpe de 0,930 impiden una afirmación concluyente de Sharpe ajustado por selección.

Los costes son constantes, 5 pb de comisión y 10 pb de slippage. Esto demuestra que el resultado no
requiere coste cero, pero no modela impacto de mercado, bid--ask dinámico, capacidad ni fiscalidad.
Dado el turnover de selección, 359\% anualizado, una especificación dependiente de liquidez puede
modificar la cartera sin alterar el Rank-IC. El universo estadounidense de gran capitalización limita
además la extrapolación a pequeñas empresas, otros países o regímenes de liquidez diferentes.

La lectura correcta es específica: multiplicidad limita el Sharpe, solapamiento limita la precisión,
cobertura limita la generalización y cartera limita la transferencia de señal. Ninguna de estas
salvedades permite ignorar los placebos, la permutación o la estabilidad por eras; cada conclusión
debe mantenerse dentro del alcance de la evidencia que realmente la respalda.
